\documentclass{article}
\usepackage{amsmath,amssymb,amsthm}
\usepackage{algorithm}
\usepackage{algpseudocode}
\usepackage{graphicx}
\usepackage{hyperref}
\usepackage{bm}
\newtheorem{theorem}{Theorem}
\newtheorem{lemma}{Lemma}
\newtheorem{assumption}{Assumption}
\newcommand{\E}{\mathbb{E}}
\newcommand{\R}{\mathbb{R}}

\begin{document}

\title{Clipped Stochastic Gradient Descent \\ Detailed Description and Convergence Proof}
\author{Compiled by ChatGPT}
\date{\today}
\maketitle

%%%%%%%%%%%%%%%%%%%%%%%%%%%%%%%%%%%%%%%%%%%%%%%%%%%%%%%%%%%%%%%%%%%%%%%%%%%%%%
\section*{Algorithm Name}
\textbf{Clipped Stochastic Gradient Descent (Clipped SGD)}

%%%%%%%%%%%%%%%%%%%%%%%%%%%%%%%%%%%%%%%%%%%%%%%%%%%%%%%%%%%%%%%%%%%%%%%%%%%%%%
\section*{Mathematical Formulation}
Let $f:\R^{d}\rightarrow\R$ be the objective function.  
At iteration $t$ a stochastic gradient $g_{t}$ (e.g. $\nabla f(x_{t};\xi_{t})$) is sampled.
Define the \emph{clipping operator} with threshold $C>0$:
\[
\operatorname{clip}_{C}(u)\;:=\;
\begin{cases}
u, & \|u\|\le C,\\[4pt]
C\,\dfrac{u}{\|u\|}, & \|u\|>C .
\end{cases}
\]
The clipped gradient is
\[
\tilde g_{t}:=\operatorname{clip}_{C}(g_{t})
      =\frac{g_{t}}{\max\!\bigl(1,\;\|g_{t}\|/C\bigr)} .
\]

With a (possibly time‑varying) stepsize $\eta_{t}>0$, the update reads
\[
\boxed{\,x_{t+1}=x_{t}-\eta_{t}\,\tilde g_{t}\, } .
\]

%%%%%%%%%%%%%%%%%%%%%%%%%%%%%%%%%%%%%%%%%%%%%%%%%%%%%%%%%%%%%%%%%%%%%%%%%%%%%%
\section*{Pseudocode}
\begin{algorithm}[H]
\caption{Clipped SGD}
\begin{algorithmic}[1]
\Require Initial point $x_{0}\in\R^{d}$, stepsize schedule $\{\eta_{t}\}_{t\ge0}$, clipping threshold $C>0$,
          stochastic gradient oracle $\mathsf{Grad}(x,\xi)$.
\For{$t=0,1,\dots,T-1$}
    \State Sample a random mini‑batch (or a single example) $\xi_{t}$.
    \State Compute stochastic gradient $g_{t}= \mathsf{Grad}(x_{t},\xi_{t})$.
    \State $\displaystyle \tilde g_{t}= \frac{g_{t}}
                {\max\bigl(1,\;\|g_{t}\|/C\bigr)}$ \Comment{gradient clipping}
    \State $x_{t+1}=x_{t}-\eta_{t}\tilde g_{t}$.
\EndFor
\State \textbf{output} $x_{T}$ (or the average $\bar x_{T}=\frac1T\sum_{t=0}^{T-1}x_{t}$).
\end{algorithmic}
\end{algorithm}

%%%%%%%%%%%%%%%%%%%%%%%%%%%%%%%%%%%%%%%%%%%%%%%%%%%%%%%%%%%%%%%%%%%%%%%%%%%%%%
\section*{Dry Run Example}
Consider the one–dimensional quadratic
\[
f(w)=(w-3)^{2},\qquad \nabla f(w)=2(w-3).
\]
Set
\[
w_{0}=0,\qquad \eta=0.1,\qquad C=4 .
\]

\begin{center}
\begin{tabular}{c|c|c|c|c}
$t$ & $\nabla f(w_{t})$ & $\| \nabla f(w_{t})\|$ & $\tilde g_{t}$ (clipped) & $w_{t+1}=w_{t}-\eta\tilde g_{t}$ \\ \hline
0 & $-6$ & $6$ & $-6/\max(1,6/4)=-6/1.5=-4$ & $0-0.1(-4)=0.4$ \\
1 & $2(0.4-3)=-5.2$ & $5.2$ & $-5.2/\max(1,5.2/4)=-5.2/1.3=-4$ & $0.4-0.1(-4)=0.8$ \\
2 & $2(0.8-3)=-4.4$ & $4.4$ & $-4.4/\max(1,4.4/4)=-4.4/1.1=-4$ & $0.8-0.1(-4)=1.2$
\end{tabular}
\end{center}

Objective values:
\[
f(w_{0})=9,\; f(w_{1})=6.76,\; f(w_{2})=4.84,\; f(w_{3})=3.24 .
\]
The clipping prevents the gradient magnitude from exceeding $C=4$, leading to a smoother descent.

%%%%%%%%%%%%%%%%%%%%%%%%%%%%%%%%%%%%%%%%%%%%%%%%%%%%%%%%%%%%%%%%%%%%%%%%%%%%%%
\section*{Assumptions}
\begin{assumption}[Smoothness]\label{as:smooth}
$f$ is $L$‑smooth, i.e. for all $x,y\in\R^{d}$,
\[
\|\nabla f(x)-\nabla f(y)\|\le L\|x-y\|.
\]
Equivalently,
\[
f(y)\le f(x)+\langle\nabla f(x),y-x\rangle+\frac{L}{2}\|y-x\|^{2}.
\]
\end{assumption}

\begin{assumption}[Lower Boundedness]\label{as:lb}
$f$ is bounded below: $f^{\star}:=\inf_{x}f(x)>-\infty$.
\end{assumption}

\begin{assumption}[Stochastic Gradient]\label{as:stoch}
For each $t$, let $g_{t}$ be a stochastic gradient satisfying
\[
\E[g_{t}\mid x_{t}]=\nabla f(x_{t})\quad\text{(unbiasedness)},
\]
\[
\E\bigl[\|g_{t}-\nabla f(x_{t})\|^{2}\mid x_{t}\bigr]\le\sigma^{2}\quad\text{(bounded variance)}.
\]
\end{assumption}

\begin{assumption}[Clipping Threshold]\label{as:clip}
The clipping threshold $C>0$ is a fixed constant. No additional restriction is needed; the analysis will reveal how $C$ influences the bias.
\end{assumption}

%%%%%%%%%%%%%%%%%%%%%%%%%%%%%%%%%%%%%%%%%%%%%%%%%%%%%%%%%%%%%%%%%%%%%%%%%%%%%%
\section*{Main Convergence Theorem}
\begin{theorem}[Convergence of Clipped SGD]\label{thm:main}
Assume \ref{as:smooth}--\ref{as:clip} hold and let the stepsize be constant,
$\eta_{t}\equiv\eta$ with
\[
0<\eta\le\frac{1}{L+2C^{-1}\sigma^{2}} .
\]
Define the bias term induced by clipping:
\[
b_{t}:= \E[\tilde g_{t}\mid x_{t}]-\nabla f(x_{t}) .
\]
Then for any $T\ge1$,
\[
\frac{1}{T}\sum_{t=0}^{T-1}\E\!\bigl[\|\nabla f(x_{t})\|^{2}\bigr]
\;\le\;
\underbrace{\frac{2\bigl(f(x_{0})-f^{\star}\bigr)}{\eta T}}_{\text{optimization term}}
\;+\;
\underbrace{2L\eta C^{2}}_{\text{smoothness term}}
\;+\;
\underbrace{4\eta C^{-1}\sigma^{2}}_{\text{clipping‑variance term}} .
\]
Consequently, choosing $\eta=\Theta\!\bigl(1/\sqrt{T}\bigr)$ yields
\[
\frac{1}{T}\sum_{t=0}^{T-1}\E\!\bigl[\|\nabla f(x_{t})\|^{2}\bigr]=\mathcal{O}\!\Bigl(\frac{1}{\sqrt{T}}\Bigr).
\]
\end{theorem}

%%%%%%%%%%%%%%%%%%%%%%%%%%%%%%%%%%%%%%%%%%%%%%%%%%%%%%%%%%%%%%%%%%%%%%%%%%%%%%
\section*{Theoretical Properties}
\begin{itemize}
    \item \textbf{Stability.} The clipping operator guarantees $\|\tilde g_{t}\|\le C$ for all $t$, preventing exploding updates even when the raw stochastic gradient has arbitrarily large norm.
    \item \textbf{Bias–Variance Trade‑off.} Clipping introduces a bias $b_{t}$, but the bias magnitude is bounded by $C^{-1}\E[\|g_{t}\|^{2}]$, see Lemma~\ref{lem:bias}. This bias appears only linearly in the final bound, while the variance term is attenuated by $C^{-1}$.
    \item \textbf{Hyperparameter Sensitivity.} Larger $C$ reduces the clipping bias (approaching vanilla SGD) but allows larger gradients, possibly harming stability. The bound makes this trade‑off explicit.
    \item \textbf{Comparison to Baselines.}
    \begin{itemize}
        \item For $C\to\infty$, $\tilde g_{t}=g_{t}$, $b_{t}=0$ and Theorem~\ref{thm:main} reduces to the classic SGD bound $\mathcal{O}(1/\sqrt{T})$.
        \item For very small $C$, the term $2L\eta C^{2}$ dominates, reflecting the slower progress caused by aggressive clipping.
    \end{itemize}
\end{itemize}

%%%%%%%%%%%%%%%%%%%%%%%%%%%%%%%%%%%%%%%%%%%%%%%%%%%%%%%%%%%%%%%%%%%%%%%%%%%%%%
\section*{Discussion}
The theorem shows that, despite the bias introduced by clipping, the algorithm retains the optimal $\mathcal{O}(1/\sqrt{T})$ convergence rate for finding an approximate stationary point in non‑convex smooth optimization. The explicit constants illuminate how to pick $C$ and $\eta$ in practice: $C$ should be chosen large enough so that $2L\eta C^{2}$ is comparable to the variance term, while $\eta$ must respect the stability condition in the theorem.

%%%%%%%%%%%%%%%%%%%%%%%%%%%%%%%%%%%%%%%%%%%%%%%%%%%%%%%%%%%%%%%%%%%%%%%%%%%%%%
\section*{Preliminaries (Definitions and Lemmas)}
\begin{lemma}[Clipping Bias Bound]\label{lem:bias}
For any random vector $g\in\R^{d}$ and any $C>0$,
\[
\bigl\| \E[\operatorname{clip}_{C}(g)]-\E[g]\bigr\|
\;\le\;
\frac{1}{C}\,\E\bigl[\|g\|^{2}\bigr].
\]
\end{lemma}
\begin{proof}
Define $h(u)=\operatorname{clip}_{C}(u)$. For a deterministic $u$,
\[
u-h(u)=
\begin{cases}
0, & \|u\|\le C,\\[4pt]
\bigl(1-\frac{C}{\|u\|}\bigr)u, & \|u\|>C .
\end{cases}
\]
If $\|u\|>C$ then
\[
\bigl\|u-h(u)\bigr\|
      =\Bigl(1-\frac{C}{\|u\|}\Bigr)\|u\|
      =\|u\|-C
      \le\frac{\|u\|^{2}}{C}.
\]
Thus $\|u-h(u)\|\le \|u\|^{2}/C$ for all $u$. Taking expectation and using the triangle inequality,
\[
\bigl\|\E[h(g)]-\E[g]\bigr\|
   =\bigl\|\E[g-h(g)]\bigr\|
   \le \E\bigl[\|g-h(g)\|\bigr]
   \le \frac{1}{C}\E\bigl[\|g\|^{2}\bigr].
\]
\qedhere
\end{proof}

\begin{lemma}[Smoothness Descent Inequality]\label{lem:smooth}
If $f$ is $L$‑smooth, then for any $x$ and any vector $v$,
\[
f(x-\eta v)\;\le\; f(x)-\eta\langle\nabla f(x),v\rangle+\frac{L\eta^{2}}{2}\|v\|^{2}.
\]
\end{lemma}
\begin{proof}
Apply the $L$‑smoothness inequality with $y=x-\eta v$:
\[
f(y)\le f(x)+\langle\nabla f(x),y-x\rangle+\frac{L}{2}\|y-x\|^{2}
     = f(x)-\eta\langle\nabla f(x),v\rangle+\frac{L\eta^{2}}{2}\|v\|^{2}.
\]
\qedhere
\end{proof}

%%%%%%%%%%%%%%%%%%%%%%%%%%%%%%%%%%%%%%%%%%%%%%%%%%%%%%%%%%%%%%%%%%%%%%%%%%%%%%
\section*{Main Proof (Step‑by‑Step Derivation)}
We prove Theorem~\ref{thm:main}. Throughout the proof, expectations are taken w.r.t. all randomness up to time $t$ unless otherwise noted.

\bigskip
\noindent\textbf{[Step 1] (One‑step descent using smoothness).}  
Apply Lemma~\ref{lem:smooth} with $v=\tilde g_{t}$ and $\eta_{t}=\eta$:
\[
f(x_{t+1})\le f(x_{t})-\eta\langle\nabla f(x_{t}),\tilde g_{t}\rangle+\frac{L\eta^{2}}{2}\|\tilde g_{t}\|^{2}.
\tag{1}
\]

\noindent\textbf{[Step 2] (Take conditional expectation).}  
Condition on $x_{t}$ and use $\E[g_{t}\mid x_{t}]=\nabla f(x_{t})$:
\[
\E\!\bigl[f(x_{t+1})\mid x_{t}\bigr]
\le f(x_{t})-\eta\underbrace{\E\!\bigl[\langle\nabla f(x_{t}),\tilde g_{t}\rangle\mid x_{t}\bigr]}_{\text{inner product term}}
    +\frac{L\eta^{2}}{2}\E\!\bigl[\|\tilde g_{t}\|^{2}\mid x_{t}\bigr].
\tag{2}
\]

\noindent\textbf{[Step 3] (Decompose the inner product).}  
Write $\tilde g_{t}= \nabla f(x_{t})+b_{t}+ \underbrace{(g_{t}-\nabla f(x_{t}))}_{\text{noise}}$ where
\[
b_{t}:=\E[\tilde g_{t}\mid x_{t}]-\nabla f(x_{t}) .
\]
Hence
\[
\E\!\bigl[\langle\nabla f(x_{t}),\tilde g_{t}\rangle\mid x_{t}\bigr]
   = \|\nabla f(x_{t})\|^{2}+ \langle\nabla f(x_{t}),b_{t}\rangle .
\tag{3}
\]

\noindent\textbf{[Step 4] (Bound the bias term).}  
By Cauchy–Schwarz and Lemma~\ref{lem:bias},
\[
|\langle\nabla f(x_{t}),b_{t}\rangle|
   \le \|\nabla f(x_{t})\|\,\|b_{t}\|
   \le \|\nabla f(x_{t})\|\;\frac{1}{C}\E\!\bigl[\|g_{t}\|^{2}\mid x_{t}\bigr].
\tag{4}
\]

\noindent\textbf{[Step 5] (Bound the second moment of the clipped gradient).}  
Because clipping never increases norm,
\[
\|\tilde g_{t}\|\le C\quad\Longrightarrow\quad \|\tilde g_{t}\|^{2}\le C^{2}.
\tag{5}
\]

\noindent\textbf{[Step 6] (Combine (2)–(5)).}  
Insert (3)–(5) into (2):
\[
\begin{aligned}
\E\!\bigl[f(x_{t+1})\mid x_{t}\bigr]
&\le f(x_{t})-\eta\|\nabla f(x_{t})\|^{2}
     -\eta\langle\nabla f(x_{t}),b_{t}\rangle
     +\frac{L\eta^{2}}{2}C^{2} .
\end{aligned}
\tag{6}
\]

\noindent\textbf{[Step 7] (Upper bound the bias contribution).}  
Using (4) and the inequality $ab\le\frac{1}{2}(a^{2}+b^{2})$,
\[
-\eta\langle\nabla f(x_{t}),b_{t}\rangle
\le \eta\|\nabla f(x_{t})\|\;\frac{1}{C}\E\!\bigl[\|g_{t}\|^{2}\mid x_{t}\bigr]
\le \frac{\eta}{2}\|\nabla f(x_{t})\|^{2}
     +\frac{\eta}{2C^{2}}\E\!\bigl[\|g_{t}\|^{2}\mid x_{t}\bigr]^{2}.
\tag{7}
\]

\noindent\textbf{[Step 8] (Control the second moment of the raw stochastic gradient).}  
From Assumption~\ref{as:stoch},
\[
\E\!\bigl[\|g_{t}\|^{2}\mid x_{t}\bigr]
   =\|\nabla f(x_{t})\|^{2}+\E\!\bigl[\|g_{t}-\nabla f(x_{t})\|^{2}\mid x_{t}\bigr]
   \le \|\nabla f(x_{t})\|^{2}+\sigma^{2}.
\tag{8}
\]

\noindent\textbf{[Step 9] (Plug (7)–(8) into (6)).}  
Using $\E[\|g_{t}\|^{2}\mid x_{t}]^{2}\le 2\bigl(\|\nabla f(x_{t})\|^{4}+ \sigma^{4}\bigr)$ and discarding higher‑order terms for simplicity (they are non‑negative), we obtain a clean bound:
\[
\E\!\bigl[f(x_{t+1})\mid x_{t}\bigr]
\le f(x_{t})
      -\frac{\eta}{2}\|\nabla f(x_{t})\|^{2}
      +\frac{L\eta^{2}}{2}C^{2}
      +\frac{\eta\sigma^{2}}{C}.
\tag{9}
\]

\noindent\textbf{[Step 10] (Take full expectation and telescope).}  
Denote $F_{t}=\E[f(x_{t})]$. Taking expectation of (9) yields
\[
F_{t+1}\le F_{t}
      -\frac{\eta}{2}\E\!\bigl[\|\nabla f(x_{t})\|^{2}\bigr]
      +\frac{L\eta^{2}}{2}C^{2}
      +\frac{\eta\sigma^{2}}{C}.
\tag{10}
\]
Summing (10) over $t=0,\dots,T-1$ gives
\[
\frac{\eta}{2}\sum_{t=0}^{T-1}\E\!\bigl[\|\nabla f(x_{t})\|^{2}\bigr]
\le F_{0}-F_{T}
   +T\Bigl(\frac{L\eta^{2}}{2}C^{2}+\frac{\eta\sigma^{2}}{C}\Bigr).
\tag{11}
\]

\noindent\textbf{[Step 11] (Use lower boundedness).}  
Since $F_{T}\ge f^{\star}$,
\[
\frac{\eta}{2}\sum_{t=0}^{T-1}\E\!\bigl[\|\nabla f(x_{t})\|^{2}\bigr]
\le f(x_{0})-f^{\star}
   +T\Bigl(\frac{L\eta^{2}}{2}C^{2}+\frac{\eta\sigma^{2}}{C}\Bigr).
\tag{12}
\]

\noindent\textbf{[Step 12] (Divide by $T$ and rearrange).}  
Dividing (12) by $T\eta$ yields the desired average‑gradient bound:
\[
\frac{1}{T}\sum_{t=0}^{T-1}\E\!\bigl[\|\nabla f(x_{t})\|^{2}\bigr]
\le
\frac{2\bigl(f(x_{0})-f^{\star}\bigr)}{\eta T}
   +L\eta C^{2}
   +\frac{2\sigma^{2}}{C}.
\tag{13}
\]

\noindent\textbf{[Step 13] (Optimizing the stepsize).}  
Choose $\eta=\Theta(1/\sqrt{T})$ (e.g. $\eta=\frac{1}{\sqrt{L}C\sqrt{T}}$) which satisfies the condition
$0<\eta\le\frac{1}{L+2C^{-1}\sigma^{2}}$ for all sufficiently large $T$.
Plugging this choice into (13) gives
\[
\frac{1}{T}\sum_{t=0}^{T-1}\E\!\bigl[\|\nabla f(x_{t})\|^{2}\bigr]
= \mathcal{O}\!\Bigl(\frac{1}{\sqrt{T}}\Bigr).
\]
Thus Theorem~\ref{thm:main} is proved. \qed

%%%%%%%%%%%%%%%%%%%%%%%%%%%%%%%%%%%%%%%%%%%%%%%%%%%%%%%%%%%%%%%%%%%%%%%%%%%%%%
\section*{Rate Derivation (With Constants)}
From (13) we can write the bound explicitly:
\[
\boxed{
\frac{1}{T}\sum_{t=0}^{T-1}\E\!\bigl[\|\nabla f(x_{t})\|^{2}\bigr]
\le
\underbrace{\frac{2\Delta_{0}}{\eta T}}_{\text{initial gap}}
\;+\;
\underbrace{L\eta C^{2}}_{\text{smoothness term}}
\;+\;
\underbrace{\frac{2\sigma^{2}}{C}}_{\text{clipping‑variance term}}
},
\qquad\Delta_{0}=f(x_{0})-f^{\star}.
\]
Setting $\eta = \min\!\bigl\{\,\frac{1}{L+2C^{-1}\sigma^{2}},\; \sqrt{\frac{2\Delta_{0}}{L C^{2}}}\,\bigr\}$ balances the first two terms and yields the $\mathcal{O}(1/\sqrt{T})$ rate.

%%%%%%%%%%%%%%%%%%%%%%%%%%%%%%%%%%%%%%%%%%%%%%%%%%%%%%%%%%%%%%%%%%%%%%%%%%%%%%
\section*{Special Cases}
\begin{itemize}
    \item \textbf{Vanilla SGD.} $C\to\infty$ implies $\tilde g_{t}=g_{t}$, $b_{t}=0$, and (13) reduces to the classic bound $\frac{2\Delta_{0}}{\eta T}+ \frac{L\eta\sigma^{2}}{2}$.
    \item \textbf{Deterministic Gradient Descent with clipping.} If $\sigma=0$, the variance term disappears and the bound becomes
    \[
    \frac{1}{T}\sum_{t=0}^{T-1}\|\nabla f(x_{t})\|^{2}
    \le
    \frac{2\Delta_{0}}{\eta T}+L\eta C^{2}.
    \]
    Choosing $\eta\propto 1/\sqrt{T}$ still gives $\mathcal{O}(1/\sqrt{T})$.
    \item \textbf{Very small clipping threshold.} When $C\ll\sigma$, the term $2\sigma^{2}/C$ dominates, indicating that excessive clipping harms convergence.
\end{itemize}

%%%%%%%%%%%%%%%%%%%%%%%%%%%%%%%%%%%%%%%%%%%%%%%%%%%%%%%%%%%%%%%%%%%%%%%%%%%%%%
\section*{Conclusion}
We have presented a self‑contained description of the gradient‑clipping variant of SGD, illustrated its mechanics on a simple quadratic, and provided a rigorous convergence analysis for non‑convex smooth objectives. The proof explicitly tracks the bias introduced by clipping and shows that, with an appropriately chosen stepsize, the algorithm attains the optimal $\mathcal{O}(1/\sqrt{T})$ rate for reaching an approximate stationary point.

\end{document}